\begin{tikzpicture}
	\tikzmath{
		\widthpic = 8;
		\small = \widthpic / 400; %the filling rectangle in the nodes should be a bit submerged
		\l = 4; %distance between the nodes
	}

	%CENTRAL NODE PARAMETERS
	\tikzmath{
		\r = 0.8;
	}

	% draw central node
	\draw[fill=cnw] (0, 0) circle (\r);

	\begin{scope}[rotate=45]
		%characterstics of 2nd node and properties
		\tikzmath{
			\d = 0.6; %thickness of the capillary
		}
		\tikzmath{
			\pb = sqrt(\r * \r - \d * \d / 4); % p border
			\pr = \pb - \small; %p length
		}
		\begin{scope}[shift={(\pr, 0)}]
			%default one
			\fill[cnw] (0, \d/2) rectangle (\l, -\d/2);
			%draw the main outer borders
			\draw[] (0, \d/2) -- (\l, \d/2);
			\draw[] (0, -\d/2) -- (\l, -\d/2);
		\end{scope}
	\end{scope}


	\begin{scope}[rotate=135]
		%characterstics of 2nd node and properties
		\tikzmath{
			\d = 0.4; %thickness of the capillary
		}
		\tikzmath{
			\pb = sqrt(\r * \r - \d * \d / 4); % p border
			\pr = \pb - \small; %p length
		}
		\begin{scope}[shift={(\pr, 0)}]
			%default one
			\fill[cnw] (0, \d/2) rectangle (\l, -\d/2);
			%draw the main outer borders
			\draw[] (0, \d/2) -- (\l, \d/2);
			\draw[] (0, -\d/2) -- (\l, -\d/2);
		\end{scope}
	\end{scope}

	%tube
	\begin{scope}[rotate=-45]
		%characterstics of 2nd node and properties
		\tikzmath{
			\d = 0.8; %thickness of the capillary
		}
		\tikzmath{
			\pb = sqrt(\r * \r - \d * \d / 4); % p border
			\pr = \pb - \small; %p length
		}
		\begin{scope}[shift={(\pr, 0)}]
			%default one
			\fill[cnw] (0, \d/2) rectangle (\l, -\d/2);
			%draw the main outer borders
			\draw[] (0, \d/2) -- (\l, \d/2);
			\draw[] (0, -\d/2) -- (\l, -\d/2);
		\end{scope}
	\end{scope}




	\begin{scope}[rotate=-135]
		%characterstics of 2nd node and properties
		\tikzmath{
			\d = 0.8; %thickness of the capillary
		}
		\tikzmath{
			\pb = sqrt(\r * \r - \d * \d / 4); % p border
			\pr = \pb - \small; %p length
		}
		\begin{scope}[shift={(\pr, 0)}]
			%default one
			\fill[cw] (0, \d/2) rectangle (\l, -\d/2);

			%draw a slob of rectangle

			\begin{scope}
				\tikzmath{
					\pos1 = 0;
					\pos2 = 0.2;
				}
				\fill[cnw] (\pos1 * \l, \d/2) rectangle (\pos2 * \l, -\d/2);
			\end{scope}

			%drawing a meniscus
			\begin{scope}
				%SET MENISCUS POSITION
				\tikzmath{
					\meniscuspos = 0.2; %in 0-1 coordinates
					\meniscusrot = 0; %0=) 180=(
				}
				%drawing meniscus
				\begin{scope}[rotate=\meniscusrot, shift={(\meniscuspos*\l - \d/4,0)}]
					\fill[cw] (0, \d / 2) rectangle (\d/2, -\d / 2);
					\fill[cnw] (0, -\d/2) arc(-90:90:\d /2) -- cycle;
					\draw[] (0, -\d/2) arc(-90:90:\d /2);
				\end{scope}
			\end{scope}

			%draw the main outer borders
			\draw[] (0, \d/2) -- (\l, \d/2);
			\draw[] (0, -\d/2) -- (\l, -\d/2);

		\end{scope}
	\end{scope}


\end{tikzpicture}
